% answers_short.tex —— 2009 年·普通高等学校招生考试·新课标卷（琼、宁卷）·理科 · 简短答案速查
% 由 tools/make_short_answers.py 自 answers.tex 的 \gksolhead/\gkanswer 提取，
% 并与 2009/2009_数学真题_新课标理_答案要点.md 逐题比对后生成，勿手改（改源 MD 后重跑）。

\begin{gknotice}
\textbf{使用说明}：本节只给最终答案便于快速核对，完整推导见解析版。答案为编写者独立推导，\textbf{非官方标准答案}。
\end{gknotice}

\gkpart{选择题}
\begin{center}\begin{tabular}{ccccccccccccc}
\toprule
题号 & 1 & 2 & 3 & 4 & 5 & 6 & 7 & 8 & 9 & 10 & 11 & 12 \\
\midrule
答案 & A & D & C & A & A & B & C & D & C & B & A & C \\
\bottomrule
\end{tabular}\end{center}

\gkpart{填空题}
\begin{itemize}[leftmargin=2.6em]
  \item[\textbf{13.}] $y=x$。
  \item[\textbf{14.}] $\frac{9\pi}{10}$。
  \item[\textbf{15.}] $140$。
  \item[\textbf{16.}] $10$。
\end{itemize}

\gkpart{解答题}
\begin{itemize}[leftmargin=2.6em]
  \item[\textbf{17.}] 方案不唯一，以下为一种可行方案。在 $A$ 处测得对 $M$、$N$ 的俯角 $\alpha$、$\gamma$，在 $B$ 处测得对 $M$、$N$ 的俯角 $\beta$、$\delta$，并测得 $AB=d$（俯角即视线与水平方向的夹角，与仰角按对顶角等价处理）。
  \item[\textbf{18.}] （1）$\frac1{100}$；（2）$x=5$，$y=15$，$B$ 类工人中个体间的差异程度更小；$A$ 类、$B$ 类以及全厂工人生产能力的平均数估计值分别为 $123$、$133.8$、$131.1$。
  \item[\textbf{19.}] （1）证明见解析；（2）$30^\circ$；（3）存在，$SE:EC=2:1$。
  \item[\textbf{20.}] （1）$\frac{x^2}{16}+\frac{y^2}{7}=1$；（2）轨迹方程为 $\lambda^2y^2=7+\left(\frac9{16}-\lambda^2\right)x^2$（$|x|\leqslant4$）；$\lambda=\frac34$ 时为两条平行线段 $y=\pm\frac{4\sqrt7}{3}$，$\lambda>\frac34$ 时为椭圆（的一部分），$0<\lambda<\frac34$ 时为双曲线（的一部分）。
  \item[\textbf{21.}] （1）单调增区间为 $(-\infty,-3)$，$(0,3)$，单调减区间为 $(-3,0)$，$(3,+\infty)$；（2）见解析。按题设条件实际推出 $\beta-\alpha>6$；原页印作「$\beta-\alpha<6$」，与条件矛盾，疑为「$>6$」之误，照录原页并登记问题条目。
\end{itemize}

\gkpart{选做题}
\begin{itemize}[leftmargin=2.6em]
  \item[\textbf{22.}] （1）（2）证明见解析。
  \item[\textbf{23.}] （1）$C_1$ 是圆 $(x+4)^2+(y-3)^2=1$，$C_2$ 是椭圆 $\frac{x^2}{64}+\frac{y^2}{9}=1$；（2）$\frac{8\sqrt5}{5}$。
  \item[\textbf{24.}] （1）$y=\begin{cases}160-10x,&0\leqslant x\leqslant10,\\ 80-2x,&10\leqslant x\leqslant20,\\ 10x-160,&20\leqslant x\leqslant30;\end{cases}$（2）$9\leqslant x\leqslant23$。
\end{itemize}
